% revised: CC 8 October 2026
% What this holds: \part{Community Input} and all four of its sections.
% Read by arlington-bsap.tex. Prose lives here; formatting lives in the wrapper and style/.
\part{Community Input}
\sectionbreaksfalse  % only Part A's sections start new pages for now; the rest is short (Sally, October 4, 2026)

% Part B is the National Civic League's. What follows is an outline, not prose:
% each item is a claim from NCL's "Arlington BSaP Outline - Community Input"
% (Drive, 30 September 2026), sorted under the report's four headings, with
% NCL's own figures. Nothing here is this report's finding until NCL writes it;
% the figures are theirs and are cited to their report when it exists.
% The order differs from Parts A and C on purpose: for the community input,
% who took part comes first, then what they said (Sally, 4 October 2026).

\textit{[Part B is the National Civic League's. What follows is their
outline of September 30, 2026, sorted under the report's four headings; every
figure in it is theirs and preliminary, and the text is theirs to write.]}

The National Civic League read five kinds of public input: public comment at
Board meetings, written correspondence to the Board, input summarized in staff
reports and Board materials, the County's Resident Satisfaction Survey, and the
County's ranked choice voting surveys. Two coders coded 218 comments from about
156 commenters for stance, focus and motivating principle, and the team ran its
own cross-tabulations of the satisfaction survey by tenure, homeownership, race,
income and gender. % waits on: comments-universe

\section{Demographic Representation}

% Composition stated once, as groundwork for everything that follows, and as a
% finding rather than a limitation (decided 30 September 2026).
\begin{itemize}
  \item Public comment is self-selected: residents with the time and the
    familiarity to take part. Both surveys over-represent White, homeowning and
    higher-income residents, which limits what they can say about renters and
    communities of color.
  \item Who is missing is itself a finding: respondents who decline to give
    their demographics are the least trusting group in the data, and Latino
    voters lag on ranked-choice awareness and understanding, so the input may
    understate both disengagement and distrust.
  % Drafting note, not an NCL finding: report every composition dimension,
  % gender beside race, homeownership, income and tenure, a null result being
  % a result (decided 30 September 2026).
  \item Familiarity and trust run opposite: long-tenured homeowners, the
    residents most familiar with the structure, report the lowest trust; newer
    residents and renters report more. Commenters are a highly engaged group,
    so the comment record may over-represent the most skeptical.
  \item The parallel in Part A: who has held the Board's seats, by the same
    groups (Section~\ref{sec:who}); the one dimension Part A cannot yet match is
    whether members rented or owned their homes.
    % waits on: coffey-renter
\end{itemize}

\section{Board Structure}

\begin{itemize}
  \item Most comments support moving forward with a review of the County's
    structure, and nearly four in five address the review process itself
    (whether, scope, how) rather than the structure on its merits.
  \item Supporters and opponents argue on different grounds: supporters on
    equity, especially racial equity, transparency and accountability;
    opponents on the process (its purpose, timing, cost, and whether renters
    are equitably included), not in defense of the current structure.
  \item Racial equity is the central principle in the comments: 71 percent of
    comments citing representational equity were supportive. The survey data
    complicate the picture, because non-White residents are under-represented
    in the satisfaction survey.
  \item Housing status is a theme across both sources: comments opposed to
    reform cite fears of renter under-representation; the satisfaction survey
    finds renters more satisfied with the current form than owners, on a thin
    sample.
  \item Transparency of County decision-making is the one satisfaction item
    that has not improved since 2022, with 35 percent satisfied in 2026, and
    comments name recent County decisions that degraded trust. Residents rate
    the County's general communication well (64 percent satisfied) and its
    decision-making poorly.
  \item NCL's own table maps comment codes to survey items: "structure is right
    for this community" and accountability and independence; trust in elected
    officials and in the County Manager; transparency and public engagement;
    representational equity and the subgroup differences.
\end{itemize}

\section{Election Method}

\begin{itemize}
  \item Voters find ranked choice voting easy to use but harder to understand,
    especially how ballots are counted: 87 percent say they understand the
    count, and about half could name the winner.
  \item Renters are more supportive of ranked choice voting than owners,
    mirroring the renter-owner gap in the satisfaction survey.
  \item Latino voters lag on awareness, understanding and support, pointing to
    a voter education need.
  \item The parallel in Part A: Section~\ref{sec:method}, the methods by which
    the Board's members have been chosen.
\end{itemize}

\section{Local Authority}

% Unchecked: NCL's outline has no theme on what Virginia law lets Arlington
% decide, but whether any comment raises Dillon's Rule, the General Assembly or
% a charter has not been searched (tracker row comments-local-authority).
Whether any of the community input bears on what Virginia law lets Arlington
decide is not yet known; this section says what the comments hold, or that
none does, when Part B is written.

% NCL's closing matter, "Considerations for Future Community Engagement" (key
% questions, and engagement approaches), belongs to Part C's Public Participation
% section by the author team's division: Part B owns methods and approaches,
% Part C owns the questions (consolidated outline, 1 October 2026).

